\chapter{Introduction}

Software security is nowadays a fundamental aspect of the development and maintenance of IT systems. The increasing complexity of applications, combined with the growing interconnection between devices and services, has significantly expanded the attack surface that can be exploited by potential malicious actors. Operating systems, browsers, drivers and distributed services are critical components whose proper functioning is closely linked to the robustness of their security properties.

In this context, proper memory management plays a central role, particularly in systems developed using low-level languages such as C and C++, which are defined as non-memory-safe. Although these languages offer high performance and control over resources, they are particularly prone to memory management errors, which can result in exploitable vulnerabilities. Memory-safe languages, in fact, have automatic mechanisms for controlling the allocation and deallocation of memory blocks. In the absence of these, while memory access is faster, an important safeguard for security is lost.

Despite decades of research and the development of numerous mitigation techniques, such vulnerabilities continue to arise even in widely used and mature software, highlighting how memory security remains an ongoing challenge in the contemporary cybersecurity landscape

A particularly relevant example of this challenge is the Linux kernel, one of the most widely deployed and scrutinized pieces of software in existence, yet still subject to memory safety vulnerabilities. Given the complexity of its subsystems and the performance constraints under which it operates, the kernel cannot rely on garbage collection or other automatic memory management strategies. Instead, it largely depends on reference counting, a technique by which each kernel object maintains a counter tracking the number of active references to it, with deallocation occurring only when the counter reaches zero. While reference counting provides a structured and efficient mechanism for controlling object lifetimes, it is not immune to subtle concurrency bugs: incorrect counter manipulations, race conditions on shared objects, or improper ordering of decrement and deallocation operations can all lead to use-after-free (UAF) vulnerabilities, in which a memory region is accessed after it has been freed and potentially reused for a different purpose.

The objective of this thesis is to investigate whether static analysis techniques can reliably detect UAF vulnerabilities arising specifically from reference counting errors in the Linux kernel, a class of bugs that, as will be discussed, has proven difficult to surface using existing automated tools. To this end, we analyze CVE-2025-21756, a real-world vulnerability affecting the vsock subsystem, as a case study: we examine why it eluded prior detection, and what properties of the bug and of the analysis tools contribute to this blind spot.
\\\\
The thesis is organized as follows:
\begin{itemize}
    \item \textbf{Background}: it covers the fundamentals of memory management: stack and heap organization, pointer semantics and reference counting mechanisms; it introduces the concept of use-after-free vulnerabilities and their relationship with reference counting errors and discusses the Linux kernel architecture with particular attention to how UAF vulnerabilities manifest in that context. The chapter concludes with an overview of existing static analysis tools for UAF detection, with a focused examination of UAFX and CountDown, the two tools that we tested, central to this work;

    \item \textbf{CVE-2025-21756}: after a general overview of the vulnerability, it describes the vsock subsystem and its internal architecture, examines the test case provided by the kernel maintainers, and analyzes the root causes of the bug and the kernel functions involved;

    \item \textbf{UAFX testing: static analysis}: it documents the experimental procedure conducted using UAFX: it covers the environment setup, the compilation of the relevant kernel modules, the extraction of LLVM bitcode and the execution of the analysis, discussing the results obtained and their implications;

    \item \textbf{CountDown testing: dynamic analysis}: TODO
\end{itemize}